% Vietnamese translation for OI Public Library
% Source-File: IOI中国国家候选队论文/2003/张云亮/张云亮.ppt
\documentclass[11pt]{article}
\usepackage{oipl-vn}

\title{Bản trình chiếu: Lựa chọn thuật toán cho bài toán}
\author{Zhang Yunliang}
\date{2003}

\begin{document}
\maketitle

\origtitle{Bàn về lựa chọn thuật toán cho bài toán}
\sourcepath{IOI China candidate team papers / 2003 / Zhang Yunliang / Zhang Yunliang.ppt}

\section{Chủ đề chính}

Bài trình bày bàn về một vấn đề rất thực tế trong thi lập trình: thuật toán
tối ưu nhất về mặt lý thuyết không phải lúc nào cũng là lựa chọn tốt nhất
trong phòng thi. Khi giới hạn dữ liệu cho phép, một lời giải đơn giản hơn,
dễ kiểm thử hơn và ít lỗi hơn có thể đáng chọn hơn một cấu trúc dữ liệu phức
tạp.

Vì vậy độ phức tạp được xét theo ba mặt:
\begin{itemize}
  \item thời gian chạy;
  \item bộ nhớ;
  \item độ phức tạp lập trình, bao gồm khả năng viết đúng, kiểm thử và sửa
        lỗi trong thời gian ngắn.
\end{itemize}
Các ví dụ trong slide đều có dạng ``thuật toán chuẩn tồn tại, nhưng ta thử
một thuật toán quá độ dễ cài đặt hơn''.

\section{Ví dụ 1: đếm số phần tử không vượt quá ngưỡng}

Dữ liệu gồm dãy
\[
  A_1,A_2,A_3,\ldots,A_n
\]
và các truy vấn
\[
  B_1,B_2,B_3,\ldots,B_m.
\]
Giới hạn trong slide là
\[
  1\le n\le10^6,\quad 1\le A_i\le10^6,\quad
  1\le m\le1000,\quad 1\le B_i\le10^6.
\]
Với mỗi \(B_i\), cần in số phần tử của dãy \(A\) không lớn hơn \(B_i\).

Thuật toán chuẩn có thể dùng cây đoạn hoặc cây Fenwick sau khi nén giá trị.
Slide chọn một cách làm đơn giản hơn, khai thác miền giá trị cố định
\([1,10^6]\) và số truy vấn nhỏ. Ta dùng:
\[
  L[1..1000000]
\]
để đếm số lần mỗi giá trị xuất hiện, và
\[
  V[1..1000]
\]
để đếm theo khối độ dài \(1000\). Nếu
\[
  (k-1)1000+1\le A_i\le k1000,
\]
thì cập nhật
\[
  L[A_i]++,\qquad V[k]++.
\]

Với một truy vấn \(B_i\) thuộc khối \(k\), đáp án là
\[
  S=\sum_{j=1}^{k-1}V[j]
    +\sum_{j=(k-1)1000+1}^{B_i}L[j].
\]
Mỗi truy vấn cần nhiều nhất \(1000+1000=2000\) phép cộng. Tổng chi phí được so
sánh trong slide như sau:
\[
\begin{array}{c|c|c}
  \text{cây đoạn} & \text{chia khối} & \text{miền giá trị}\\ \hline
  (n+m)\log_2 G & 2n+2000m & G=10^6
\end{array}
\]
Với \(m\le1000\), thuật toán chia khối rất dễ cài đặt và đủ nhanh. Đồng thời
nó gợi ý tự nhiên tới các cấu trúc phân cấp hơn như chia \(100\cdot100\cdot
100\), rồi sau đó tới cây đoạn.

\section{Ví dụ 2: thống kê doanh thu}

Ví dụ thứ hai là bài ``thống kê doanh thu''. Với doanh thu ngày \(i\) là
\(a_i\), định nghĩa
\[
  d_i =
  \begin{cases}
    a_1, & i=1,\\
    \min\limits_{1\le j<i}|a_i-a_j|, & i>1.
  \end{cases}
\]
Cần tính
\[
  S=\sum_i d_i.
\]
Thuật toán chuẩn là cây cân bằng: khi đọc \(a_i\), tìm phần tử đã xuất hiện
gần nhất ở hai phía.

Slide đưa ra cách làm dễ viết hơn. Trước hết đọc toàn bộ dữ liệu, sắp xếp để
được dãy
\[
  B_1,B_2,\ldots,B_n.
\]
Mỗi giá trị \(P\) trong dữ liệu gốc được đổi thành vị trí \(q\) sao cho
\(B_q=P\). Sau đó bài toán trở thành: trong tập các vị trí đã xuất hiện, tìm
vị trí gần \(q\) nhất ở bên trái và bên phải.

Ta lại dùng mảng đếm và mảng khối, lần này với kích thước được slide ghi là
\[
  L[1..32767],\qquad V[1..327],
\]
tức mỗi khối dài \(100\). Nếu \(q\) thuộc khối \(k\), trước hết tìm trong
khối đó hai vị trí
\[
  g_a\ge q,\quad g_b\le q,\qquad L[g_a]>0,\ L[g_b]>0.
\]
Nếu một phía không tồn tại, lần lượt nhìn sang các khối lân cận có
\(V[t]>0\), rồi quét trong khối ấy để lấy phần tử gần nhất. Cuối cùng so sánh
\(|P-B_{g_a}|\) và \(|P-B_{g_b}|\) nếu các vị trí này tồn tại, cộng giá trị
nhỏ hơn vào \(S\), rồi đánh dấu \(q\) đã xuất hiện.

Các chương trình \texttt{MYTURN.PAS} và \texttt{TURNOVER.PAS} trong nguồn gốc
là các bản cài đặt của ý tưởng này và lời giải chuẩn. Theo ngoại lệ về mẫu mã,
bản dịch không chép lại toàn bộ mã Pascal; điều quan trọng của slide là cách
thay cây cân bằng bằng tiền xử lý sắp xếp cộng chia khối.

\section{So sánh và chọn lựa}

Ở cả hai ví dụ, thuật toán thay thế không phải ``đẹp'' nhất về lý thuyết. Nó
phụ thuộc vào giới hạn cụ thể: miền giá trị không quá lớn, số truy vấn nhỏ,
hoặc có thể đọc trước toàn bộ dữ liệu. Nhưng chính các giả thiết ấy làm cho
lời giải đơn giản trở nên đáng dùng.

Điểm cần cân nhắc khi chọn thuật toán:
\begin{itemize}
  \item Nếu cả hai thuật toán đều chắc chắn đạt thời gian, chọn thuật toán dễ
        viết và dễ kiểm thử hơn.
  \item Nếu thuật toán đơn giản nằm sát giới hạn thời gian, cần đánh giá
        trường hợp xấu nhất bằng số phép toán cụ thể, không chỉ bằng ký hiệu
        \(O(\cdot)\).
  \item Nếu thuật toán đơn giản là ``thuật toán quá độ'', nó vẫn có giá trị
        huấn luyện vì giúp ta nhìn ra thuật toán chuẩn ở mức trừu tượng hơn.
\end{itemize}

\section{Tổng kết}

Thông điệp của bộ slide không phải là bỏ qua thuật toán tốt. Ngược lại, khi
hiểu nhiều mức lời giải cho cùng một bài, ta biết lúc nào cần cây đoạn, cây
cân bằng hoặc cấu trúc dữ liệu mạnh hơn, và lúc nào một mảng đếm kèm chia
khối đã đủ. Trong phòng thi, lựa chọn đúng là lời giải đạt điểm trong giới
hạn thời gian với rủi ro cài đặt thấp nhất.

\end{document}
